\section{Desafios de Interoperabilidade e Governança na Arquitetura de DTs}

A materialização clínica de todo o potencial in-silico exposto defronta-se irremediavelmente com as barreiras estruturais da própria indústria de tecnologia odontológica: a crise da fragmentação e da interoperabilidade.

Para que um gêmeo digital alcance o seu ápice de utilidade — operando como um modelo \textit{composicional} modular — é absolutamente mandatório que a coroa desenhada por um laboratório de prótese remoto e exportada em formato geométrico interaja perfeitamente com a simulação elástica da fisiologia periodontal calculada no supercomputador da clínica. Isto prescreve a adoção militante e estrita de \textbf{ontologias e vocabulários médicos universais} (SNOMED CT, LOINC) aliados a barramentos de mensageria (\textit{message brokers}) e estruturas abertas de representação (HL7 FHIR, DICOM, STL/PLY sem encriptação obscura).

O ecossistema odontológico tradicionalmente apoia-se em "silos" de dados, mantidos por corporações através de cadeias de arquivos em formatos restritos (\textit{vendor lock-in}) que aprisionam o dado radiológico ou oclusal ao \textit{software} proprietário do fabricante da máquina. Tais práticas predatórias ameaçam frontalmente a consolidação das simulações complexas, que precisam da confluência dos dados~\cite{robles2023opentwins}.

\destaque{A sobrevivência dos Gêmeos Digitais na saúde não está restrita a avanços em Deep Learning, mas à aceitação e construção de pipelines de dados em \textit{Open-Source}, assegurando que os dados do paciente não sejam sequestrados por formatos corporativos indecifráveis.}

Organizações de fomento acadêmico open-source, como o \textit{Digital Twin Consortium} (DTC) e a rede \textit{OpenTwins}, desempenham um papel messiânico na demolição desses silos, forjando repositórios modulares capazes de integrar de forma agnóstica desde uma varredura tomográfica de rotina até os logs metabólicos do paciente em tempo real.
